\begin{table}[H]
\centering\small
\caption{The US-demand effect by the timing of the partner-growth shift. The dependent variable is recovery duration in months; each cell is the US-demand coefficient with its country-clustered $t$-statistic. The onset measure uses only the crisis-year value of partner growth and is largely predetermined relative to the recovery, so it mitigates the window-overlap concern; the pre-sample measure uses partner growth over the five years before the downturn.}
\label{tab:prewindow}
\begin{tabularx}{\textwidth}{@{}lYY@{}}
\toprule
Timing of partner growth & Full cycle & Climb \\ \midrule
Baseline: peak to peak$+2$ & $-5.3^{***}$ {\scriptsize$(-3.5)$} & $-3.4^{***}$ {\scriptsize$(-3.0)$} \\
Onset: crisis (peak) year only & $-3.8^{***}$ {\scriptsize$(-4.1)$} & $-2.7^{***}$ {\scriptsize$(-3.9)$} \\
Pre-sample: peak$-5$ to peak$-1$ & $-2.6$ {\scriptsize$(-1.5)$} & $-1.5$ {\scriptsize$(-1.1)$} \\
\midrule
Episodes & 263 & 263 \\
Clusters & 34 & 34 \\
\bottomrule
\end{tabularx}
\\[2pt]\footnotesize{Weights are pre-crisis export shares throughout; only the years over which partner growth is averaged differ. Controls: trade openness, an indicator for an active IMF-supported arrangement, and shock depth. The US pull enters on its own here, rather than alongside the other three partner pulls as in the horse race of Table~\ref{tab:main}, so the baseline coefficient ($-5.3$) is close to but not identical to the full-cycle coefficient there ($-5.6$). $^{*}$, $^{**}$, $^{***}$: significance at 10, 5, 1 percent.}
\end{table}
